\section*{Anexo E. Comparación con otros trabajos}
\phantomsection
\addcontentsline{toc}{subsection}{Anexo E. Comparación con otros trabajos}
\makeatletter\def\@currentlabel{E}\makeatother
\label{anx:trabajos}
\setcounter{table}{0}
\renewcommand{\thetable}{E.\arabic{table}}
\renewcommand{\theHtable}{E.\arabic{table}}
% fuente de todo el anexo: fichas de lit_todas.json del 27.09.2026 (64 fichas, 54 obras), solo las de
% verificacion 'confirmado' o 'corregido'. Cada cifra se leyo en la tabla o el texto de la obra citada que
% indica su ficha; los porcentajes estan pasados a la convencion de la tesis (negativo = menos retraso,
% espera o tiempo de viaje). Las claves estan en referencias.bib desde el 27.09.2026.

Este anexo reúne los trabajos de control semafórico con los que la
Sección~\ref{sec:comparacion-literatura} compara los resultados del plan
de pruebas. La búsqueda se hizo en setiembre de 2026 y se organizó en
cuatro frentes: bancos de prueba y revisiones, control de un cruce
aislado o de un corredor, redes con varios agentes, y trabajos de
América Latina junto con sistemas que operan en ciudades. De cada obra
se registraron la red, el simulador, la demanda, el algoritmo, las
entradas y salidas del agente, la referencia contra la que se mide, si
esa referencia estaba ajustada y las cifras con su métrica. Cada cifra
se releyó en la fuente, en el texto completo cuando estuvo disponible,
y cuando un trabajo reporta varios escenarios la mediana se recalculó
con sus tablas. En once obras la relectura corrigió algún dato, como
los autores, la sede, la métrica o la referencia de una cifra. Una
investigación publicada como tesis y como artículos cuenta como una sola
obra: las cifras de MARLIN-ATSC salen de la tesis
\citep{ElTantawy2012_Tesis}, cuyos capítulos se publicaron también por
separado \citep{ElTantawy2012_ITSC,ElTantawy2014_Parametros,ElTantawy2013_MARLIN}.

Así se reunieron 54 obras. Quedaron fuera cuatro cuya cifra solo pudo
leerse en el resumen o en una fuente secundaria, entre ellas un
Q-learning para un cruce aislado \citep{Abdulhai2003_QLearning} y la
tesis de la UAQ revisada en el Capítulo~\ref{cap:estado-arte}
\citep{UAQ2024}, cuyo capítulo de resultados no estuvo accesible.
También quedó fuera un trabajo que mide el aprendizaje multiagente en un
simulador propio muy simplificado y contra un fijo de 5~s de verde
\citep{Mahato2025_MARLvsFijo}. Tampoco se usa el 73\,\% de mejora frente
a tiempo fijo que RESCO \citep{Ault2021_RESCO} atribuye a Mousavi et al.: en la fuente, esa
cifra es la reducción de la cola frente a otra red neuronal
\citep{Mousavi2017_PGDQN}. De las 49 obras restantes, siete no hacen una
comparación propia (dos simuladores o entornos y cinco revisiones) y 42
entran en las Tablas~\ref{tab:trabajos-comparados}
a~\ref{tab:trabajos-comparados-4}, una por frente y en orden
cronológico dentro de cada una.
% fuente: fichas con verificacion 'parcial' (fuera): Abdulhai2003_QLearning, Arel2010_MARL5 (sin clave
% en el .bib), CarrascoMorales2024_UAQ (clave UAQ2024) y Zamora2026_Peru (sin clave, sin texto completo).
% Mahato2025_MARLvsFijo: 'confirmado', pero con simulador en Pygame y fijo de 5 s (advertencia de su
% ficha). Mousavi: comparacion CMP4 de la ficha Ault2021_RESCO (corregida). Sin comparacion propia:
% Alegre2019_SUMORL, Zhang2019_CityFlow, Wei2019_Survey, Haydari2022_Survey, Noaeen2022_SLR,
% Yau2017_Survey y Greguric2020_Review. Obras con alguna ficha 'corregido': RESCO, LibSignal, Lioris,
% MARLIN-ATSC, Liang, Trinity, SURTRAC, Vivacity, SCOOT, FHWA (2008) y la ola verde de Lima.

Todas las cifras se expresan como cambio porcentual frente a la
referencia, con la convención de esta tesis: un valor negativo indica
menos retraso, espera o tiempo de viaje que la referencia, aunque el
trabajo original reporte la misma mejora como un porcentaje positivo.
Las cifras solo son comparables entre filas con la misma métrica. El
tiempo de viaje incluye el recorrido a velocidad libre, que ningún
control puede acortar, y por eso diluye el porcentaje. La espera
detenida es solo una parte del retraso y lo infla. En los resultados de
X-Light, el mismo IPPO reduce el tiempo de viaje un 18.9\,\% y el
retraso un 40.4\,\% en la misma malla y frente al mismo fijo
\citep{Jiang2024_XLight}; en el piloto de
SURTRAC, el tiempo de viaje baja un 25.8\,\% y la espera un 40.6\,\%
\citep{Smith2013_Surtrac}. La métrica puede incluso invertir el signo:
en la misma malla y frente a Greedy, FMA2C reduce un 90.5\,\% el retraso
medido en el cruce y aumenta un 10.8\,\% el retraso del viaje
\citep{Ma2020_FMA2C}.
% fuente: X-Light CMP3 (IPPO en Grid4x4, 94.64 a 56.38 s de retraso; -18.9 % en tiempo de viaje);
% SURTRAC CMP0 y CMP1 (global, ponderado por volumen); FMA2C CMP2 y CMP3 (malla 4x4 con flujo creciente)

En las tablas de trabajos comparados, la columna Red da el número de
intersecciones semaforizadas de las redes evaluadas y si son reales o
sintéticas, y la columna Sim., el simulador; \emph{campo} indica una
medición en calles. En la referencia, un \emph{fijo sin tunear} usa
verdes iguales o arbitrarios, un \emph{fijo no declarado} es uno cuyo
ajuste el trabajo no describe y un \emph{fijo tuneado} se ajustó por
barrido o por programación lineal. El \emph{plan real} es el programa
que opera en la ciudad o el que trae el escenario, sin reoptimizar, y
\emph{otro RL} indica que la referencia es otro aprendiz. Varias líneas
con la misma referencia son escenarios distintos. Una cifra sin
aprendiz nombrado es la del mejor aprendiz que el trabajo reporta en ese
escenario, o la mediana de sus escenarios; con nombre, es la de ese
aprendiz. \emph{Máx.} marca una cifra que el trabajo da como máximo
(``hasta''), y \emph{sin cifra}, una comparación que el trabajo solo
muestra en figuras o describe sin números.

@@COMPARADOS_1@@
% fuente: fichas de cada obra en lit_todas.json. LibSignal: mediana de 8 escenarios en SUMO (Tablas 4 y 7);
% Lioris: mediana de 7 pares O-D; Efficient Pressure: mediana de 5; Advanced-XLight: mediana de 6; T-REX:
% mediana de 3 redes reales (plan real) y de 4 escenarios (Max-Pressure), grilla 4x4 para el fijo sin
% tunear; RESCO: validacion en las mallas 4x4 (Fig. 4), IPPO frente a IDQN es la mediana de 6 tareas.

@@COMPARADOS_2@@
% fuente: MARLIN-ATSC, Tablas 5-8 (perfil uniforme), 6-12 (demanda actual) y 7-5 de la tesis; IntelliLight,
% datos reales de Jinan; FRAP, mediana de 18 cruces; PressLight, mediana de 10 escenarios.

@@COMPARADOS_3@@
% fuente: CoLight frente a Max-Pressure, mediana de los 5 escenarios de la ficha (-6.6, -12.7, -10.7,
% -29.6, -19.4), calculada para este anexo; MPLight, mediana de 5; X-Light y CoSLight frente al fijo,
% mediana de 5 escenarios; IPPO y MAPPO de CoSLight, red de Nanshan; FMA2C y MA2C, Monaco salvo la malla.

@@COMPARADOS_4@@
% fuente: SCOOT, TAL 7/99 (estimaciones promedio y Worcester); FHWA (2008), proyeccion de NTOC (rango
% de 15 a 40 %); SURTRAC, Tabla 1 (global); Higuera, Tabla 8-1 (Q-BR); Montoya, Tabla 4.2 (nivel alto);
% Tarquino, Tabla 6 (escenario 1); Uniandes, DQN y Q-learning (resumen y Sec. 8).

Las Tablas~\ref{tab:trabajos-entradas} a~\ref{tab:trabajos-entradas-3}
describen, para las obras con aprendizaje por refuerzo que publican su
formulación, qué observa el agente, qué decide, con qué recompensa
aprende y con qué métricas se evalúa. No incluyen los controles sin
aprendizaje ni el informe de Vivacity \citep{Velasco2021_Vivacity}, que
no describe su formulación. La penúltima fila resume una revisión
sistemática de 160 artículos sobre redes de dos o más cruces
\citep{Noaeen2022_SLR}, y la última describe esta tesis.

@@ENTRADAS_1@@

@@ENTRADAS_2@@

@@ENTRADAS_3@@
% fuente: campos estado, accion, recompensa y metricas de cada ficha; fila de la revision: Noaeen2022_SLR
% (preprint de engrXiv); fila de esta tesis: Capitulo III (sec:formulacion, sec:u4, Tabla tab:obs-u4 y sec:u3).

La mejora reportada depende sobre todo de la referencia. Con una cifra
por obra y por tipo de referencia, solo de trabajos con aprendizaje por
refuerzo, sin informes de fabricantes y solo en retraso, espera o
tiempo de viaje, la mediana es de $-46.0$\,\% frente a un fijo sin
tunear o no declarado (15 obras), de $-35.7$\,\% frente a un fijo
tuneado, un plan de Webster o el plan real (10 obras), de $-20.7$\,\%
frente al actuado (4 obras) y de $-15.5$\,\% frente a Max-Pressure u
otra heurística adaptativa (17 obras). Buena parte de la distancia
entre la primera categoría y la última la cubre una heurística sin
entrenamiento: frente al mismo fijo sin tunear, Max-Pressure ya reduce
el tiempo de viaje un 35.8\,\% en un trabajo \citep{Wu2021_EfficientMP}
y un 48.5\,\% en otro \citep{Mei2024_LibSignal}. Las cifras más bajas
aparecen donde la referencia está mejor ajustada. Frente a un fijo de
verde uniforme tuneado por barrido y evaluado con 32 semillas, el mejor
aprendiz reduce el tiempo de viaje un 8.9\,\%
\citep{Genders2019_Framework}. Frente al plan real de las redes de
Colonia e Ingolstadt, IPPO aumenta el tiempo de viaje en dos de tres
redes \citep{Nguyen2025_TREX}, y ningún algoritmo de EPyMARL supera al
plan en el corredor de tres cruces \citep{Ault2022_EPyMARL}.
% fuente de las medianas: calculadas para este anexo sobre las fichas (script de trabajo anexoe/gen.py de
% la sesion del 27.09.2026; se rehacen a mano con esta lista). Reglas: solo obras con RL y sin informes de
% fabricante (fuera Vivacity, Green Light y City Brain); una cifra por obra y por tipo de referencia, que es
% la mediana de sus lineas de ese tipo en las tablas (cada linea ya es la mediana de escenarios que da la
% ficha, si la da). No entran las lineas que muestran un aprendiz distinto del mejor (IPPO en LibSignal,
% T-REX, X-Light y CoSLight; DQN en Genders y Razavi 2019; Q-learning en Uniandes; PPO en Ault et al.
% 2020), las de otro RL ni las de retraso en el cruce; la linea de T-REX con IDQN en la malla si entra, porque
% es la unica de ese tipo en su ficha. Frente a una heuristica, Max-Pressure si la obra lo compara y si no
% Greedy, LQF, SOTL o SURTRAC; si la obra ajusta la heuristica, la version ajustada.
% Fijo sin tunear o no declarado (n = 15, mediana -46.0, rango -86 a -25.2): LibSignal -55.1; T-REX -28.7
%   (grilla, IDQN); LIT -60.65 (mediana de -55.5 y -65.8); Cabrejas-Egea -85.2; Efficient Pressure -38.2;
%   Advanced-XLight -38.6; Gao -86 (max.); Liang -25.7; PressLight -52.1; CoLight -25.2; MPLight -47.5;
%   AttendLight -46; X-Light -28.8; CoSLight -27.3; Uniandes -63.49. Sin cifra: Mousavi, Trinity.
% Fijo tuneado, Webster o plan real (n = 10, mediana -35.65, rango -62.49 a -8.9): Genders y Razavi 2019
%   -8.9; LIT -35.8 (mediana de -34.0 y -37.6); T-REX -17.7; MARLIN-ATSC -35.5 (mediana de -35.5, -21.4 y
%   -37.6); IntelliLight -40.3; FRAP -61.1 (mediana de -64.1 y -58.1); PressLight -50.6 (onda verde);
%   Wang -21.2; Higuera -62.49; Borges -28.4. Sin cifra: RESCO, EPyMARL.
% Actuado (n = 4, mediana -20.65, rango -48.0 a -10.1): Cabrejas-Egea -48.0 (mediana de -21.7, -48.0 y
%   -56.9); MARLIN-ATSC -21.9 (mediana de -25.1 y -18.7); Wang -10.1; Ault et al. 2020 -19.4 (max.).
% Max-Pressure u otra heuristica (n = 17, mediana -15.5, rango -47 a +22.0): RESCO -15.5 (mediana de -11 y
%   -20); Genders y Razavi 2019 +22.0; Cabrejas-Egea -26.6 (SURTRAC); Efficient Pressure -3.7; LibSignal
%   -5.7 (SUMO); Advanced-XLight -8.6; T-REX -20.2; Gao -47 (LQF, max.); IntelliLight -22.1 (SOTL);
%   PressLight -42.3; S2R2L -19.9 (ajustado); CoLight -12.7; MA2C -20.8 (Greedy); MPLight -14.4; FMA2C
%   +2.25 (mediana de -6.3 y +10.8); AttendLight -39; CoSLight -9.75 (mediana de -13.8 y -5.7).
% Max-Pressure sin RL frente al fijo sin tunear: Efficient Pressure CMP3 (-35.8, mediana de 5) y LibSignal
% CMP5 (-48.5, mediana de 8 en SUMO). IPPO frente al plan real en T-REX: +50.2, -3.3 y +5.1 (ficha
% Nguyen2025_TREX, CMP2). EPyMARL: Cologne Corridor, sin cifras en el trabajo.
